% ======================================================================
% APÉNDICE A: PROPIEDAD INTELECTUAL Y LICENCIAS
% ======================================================================

\chapter{Propiedad Intelectual y Licencias}
\label{ap:licencias}

\section{Marco Normativo de la Universidad de Granada}
\label{sec:ap_normativa_ugr}

La titularidad de los derechos de propiedad intelectual e industrial derivados de los Trabajos de Fin de Grado en la Universidad de Granada se rige por la \textit{Normativa sobre los Derechos de Propiedad Industrial e Intelectual derivados de la actividad investigadora de la Universidad de Granada} (aprobada en Consejo de Gobierno de 31 de enero de 2017, NCG115/1) y el \textit{Reglamento de Trabajo y Proyecto Fin de Grado de la Universidad de Granada} (NCG187/2).

En particular, el Artículo 14 del Reglamento de TFG establece que:
\begin{enumerate}
    \item El estudiantado tiene derecho al reconocimiento moral inalienable de la autoría de su Trabajo de Fin de Grado y a la protección de la propiedad intelectual asociada al mismo.
    \item La titularidad patrimonial de los resultados derivados del TFG corresponde al estudiante generador de los mismos, salvo en aquellos supuestos en los que el trabajo se desarrolle en el marco de proyectos de investigación financiados o convenios específicos con empresas, en cuyo caso rige el régimen de cotitularidad con la Universidad de Granada.
\end{enumerate}

\section{Licenciamiento de la Memoria y Documentación}
\label{sec:ap_licencia_memoria}

La presente memoria académica, incluyendo su texto, tablas, diagramas e interpretaciones analíticas, se publica bajo la licencia internacional \textbf{Creative Commons Reconocimiento - NoComercial - CompartirIgual 4.0 (CC BY-NC-SA 4.0)}:
\begin{itemize}
    \item \textbf{Atribución (BY):} Se debe otorgar el crédito correspondiente al autor (\emph{Emilio Gullón}), proveer un enlace a la licencia e indicar si se han realizado modificaciones.
    \item \textbf{No Comercial (NC):} No se autoriza el uso de la memoria ni de sus contenidos con fines explícitamente comerciales o mercantiles.
    \item \textbf{Compartir Igual (SA):} Cualquier obra derivada construida a partir de este documento debe difundirse obligatoriamente bajo los mismos términos de licenciamiento.
\end{itemize}

El documento íntegro queda depositado en acceso abierto en el repositorio institucional \textbf{Digibug} de la Biblioteca de la Universidad de Granada (\url{https://digibug.ugr.es}).

\section{Licenciamiento del Código Fuente (\code{fatigueset-lib})}
\label{sec:ap_licencia_software}

El código fuente de la librería modular \code{fatigueset-lib}, desarrollado durante el TFG y alojado en el repositorio público de GitHub del proyecto, se distribuye bajo la licencia de software libre permisiva \textbf{MIT License}:

\begin{lstlisting}[style=bashstyle, caption={Texto oficial de la licencia MIT de \code{fatigueset-lib}.}]
MIT License

Copyright (c) 2026 Emilio Gullon

Permission is hereby granted, free of charge, to any person obtaining a copy
of this software and associated documentation files (the "Software"), to deal
in the Software without restriction, including without limitation the rights
to use, copy, modify, merge, publish, distribute, sublicense, and/or sell
copies of the Software, and to permit persons to whom the Software is
furnished to do so, subject to the following conditions:

The above copyright notice and this permission notice shall be included in all
copies or substantial portions of the Software.

THE SOFTWARE IS PROVIDED "AS IS", WITHOUT WARRANTY OF ANY KIND, EXPRESS OR
IMPLIED, INCLUDING BUT NOT LIMITED TO THE WARRANTIES OF MERCHANTABILITY,
FITNESS FOR A PARTICULAR PURPOSE AND NONINFRINGEMENT. IN NO EVENT SHALL THE
AUTHORS OR COPYRIGHT HOLDERS BE LIABLE FOR ANY CLAIM, DAMAGES OR OTHER
LIABILITY, WHETHER IN AN ACTION OF CONTRACT, TORT OR OTHERWISE, ARISING FROM,
OUT OF OR IN CONNECTION WITH THE SOFTWARE OR THE USE OR OTHER DEALINGS IN THE
SOFTWARE.
\end{lstlisting}

Esta elección maximiza la reutilización del framework por parte de la comunidad investigadora internacional, facilitando el benchmarking y la extensión a nuevos conjuntos de datos biomédicos sin restricciones propietarias.
